\bigskip
\bigskip
\subsubsection*{Խնդիրներ և հարցեր թեմա 10-ի վերաբերյալ}

\begin{enumerate}[label=\thesection.\arabic*.]

% 10.1
\item Նկարագրեք բոլոր զուգամետ հաջորդականությունները կապակցված կե\-տա\-զույգ տարածությունում՝ $X = \{x_1, x_2\}$, $\tau = \left\{\varnothing, \{x_1\}, \{x_1, x_2\}\right\}$։

% 10.2
\item Դիցուք $X$-ը անվերջ բազմություն է վերջավոր լրացումների տոպոլոգիայով։ Ապացուցեք, որ ցանկացած $\{x_n\} \subset X$ հաջորդականություն՝ կազմված $X$-ի տարբեր կետերից, զուգամիտում է $X$-ի ցանկացած կետի։

% 10.3
\item Մաթեմատիկական անալիզի դասընթացում ապացուցվում է՝ $(\R, \text{սովոր.})$-ում մոնոտոն աճող, վերևից սահմանափակ ամեն մի հաջորդականություն ունի միակ սահման։ Ապացուցեք, որ այսպիսի հաջորդականությունները զուգամետ չեն $(\R, \mapsto)$ տարածությունում։

% 10.4
\item Դիտարկենք որևէ հաշվելի $X = \{x_1, x_2, \dots\}$ բազմություն։ Սահմանեք $X$-ի համար այնպիսի $\tau$ տոպոլոգիա, որ $(X, \tau)$ տարածությունում բոլոր զուգամետ հաջորդականություններն ունենան միևնույն սահմանը։

\begin{hint}
Քննարկեք $\left\{\varnothing, X, \{x_1\}, \{x_1, x_2\}, \{x_1, x_2, x_3\}, \dots\right\}$ ընտանիքը։
\end{hint}

% 10.5
\item Դիտարկենք որևէ $(X, \text{հաշվ․ լր․})$ տարածություն, և դիցուք $A$-ն $X$-ի որևէ ոչ հաշվելի ենթաբազմություն է։ Ապացուցեք՝
\begin{enumerate}
    \item[ա)] ցանկացած $b \in X \setminus A$ կետ հպման կետ է $A$-ի համար,
    \item[բ)] գոյություն չունի $\{a_n\} \subset A$ հաջորդականություն, որ $\lim a_n = b$։
\end{enumerate}

% 10.6
\item Ապացուցեք․ եթե $\{x_n\} \subset X$ հաջորդականությունն ունի $\lim x_n = a$ սահման, ապա $a$-ն պատկանում է $\widebar{\{x_n\}}$ փակմանը։

% 10.7
\item Ապացուցեք․ ցանկացած $\{x_n\} \subset X$ զուգամետ հաջորդականության ամեն մի $\{y_n\}$ ենթահաջորդականություն ունի նույն սահմանը (սահմանները), ինչը որ ունի $\{x_n\}$-ը։

% 10.8
\item Հաջորդականության ենթահաջորդականություն սահմանելիս պահանջեցինք $h: \mathbb{N} \to \mathbb{N}$ արտապատկերման գոյություն, որը պետք է բավարարի $h(i) > h(j)$ պայմանին ամեն անգամ, երբ $i > j$։
\par \textbf{Հարց։} Ինչպիսի՞ էֆեկտներ կարող են առաջանալ, եթե սահմանման մեջ վերո\-հիշյալ պայմանը փոխարինենք $h(i) \ge h(j)$, երբ $i > j$ ավելի թույլ պայմանով։
\par Հաջորդ երեք խնդիրներում հստակեցվում է հարցադրումը։

% 10.9
\item Ճի՞շտ է արդյոք, որ ցանկացած ստացիոնար հաջորդականության ամեն մի ենթահաջորդականություն նույնպես ստացիոնար հաջորդականություն է։

% 10.10
\item Կմնա՞ արդյոք ճիշտ պնդումը․ $\{x_n\} \subset X$ զուգամետ հաջորդականության ցան\-կացած ենթահաջորդականություն
\begin{enumerate}
    \item[ա)] նույնպես զուգամետ է,
    \item[բ)] ունի նույն սահմանները, ինչը որ ունի $\{x_n\}$ հաջորդականությունը։
\end{enumerate}

% 10.11
\item Նշանակենք $\{\lim x_n\}$ սիմվոլով $\{x_n\} \subset X$ հաջորդականության բոլոր սահ\-ման\-նե\-րի բազմությունը (այն կարող է լինել նաև $\varnothing$ բազմություն)։ Ճի՞շտ է արդյոք պնդումը․ $\{x_n\}$ հաջորդականության ցանկացած $\{y_n\}$ ենթա\-հա\-ջոր\-դա\-կա\-նու\-թյան դեպքում տեղի ունի $\{\lim y_n\} \subset \widebar{\{x_n\}}$ ներդրում։

\end{enumerate}